% Nouveau EGA V (cote 25): archive cote 25, site page 91 (PDF page 92).
% typescript. Draft compared: none.
% AI-transcribed by Claude Opus 5.5 (claude-opus-5-5, adaptive thinking, effort high), Message Batches API, batch ega25b (msgbatch_01A7n5YuXubUPNK1zsepi8RG), 2026-11-13; checked against both supplied scan parts; see CHECK_DECISIONS.tsv.
\scanpage{25/91}
\mspage{85}

Bien entendu, pour justifier dans le cas \var{\xout{actuel}}{}{typewriter strike-out}{ega25_25_091-v0001} du présent n° le passage
de \hand{$\eta$} générique aux points voisins de $\mathrm{Grass}^{m+1}(P)$, et pouvoir également
le cas échéant reprendre à notre compte les explications générales de
7.1. , il y a lieu de considérer les diagrammes
\diagram{\[\begin{array}{ccc}
X & \hand{\longleftarrow} & \overset{\hand{\sim}}{X}(C) \\
\hand{\downarrow} & & \hand{\downarrow} \\
S & \hand{\longleftarrow} & \hand{Q}(C)
\end{array}\]}
obtenus à l'aide des divers $C \hand{\in} \mathrm{Grass}^{m+1}(S)$, et \var{\inserted[typed]{plus généralement}}{plus généralement}{insertion}{ega25_25_091-v0002} ceux obtenus de même
après changement de base $T \hand{\to} S$ pour des points \hand{$\unsure{\xi|\zeta}$} \hand{$\in$} $\mathrm{Grass}^{m+1}($\var{\over{P}{\unsure{T}}}{\unsure{T}}{writing over}{ega25_25_091-v0003}$)$ :
\diagram{\[\begin{array}{ccc}
X_T & \hand{\longleftarrow} & \overset{\hand{\sim}}{X}(C_{\hand{\unsure{\xi|\zeta}}}) = \overset{\hand{\sim}}{X}_T(C_{\hand{\unsure{\xi|\zeta}}}) \\
\hand{\downarrow} & & \hand{\downarrow} \\
T & \hand{\longleftarrow} & \hand{Q}(C_{\hand{\unsure{\xi|\zeta}}})
\end{array}\]}
comme déduits, par changement\hand{s} de base \hand{$\unsure{\xi|\zeta}$} : $T \hand{\longrightarrow} \mathrm{Grass}^{m+1}$\var{\deleted{$(P)$}}{}{deletion}{ega25_25_091-v0004} $= \underline{T}$ ,
du diagramme \uline{universel} (relatif au point canonique de $\mathrm{Grass}^{m+1}$ dans $\underline{T}$ ):
\diagram{\[\begin{array}{ccc}
X_{\underline{T}} & \hand{\longleftarrow} & \overset{\hand{\sim}}{X}(\underline{C}) \\
\hand{\downarrow} & & \hand{\downarrow} \\
\underline{T} & \hand{\longleftarrow} & \hand{Q}(\underline{C})
\end{array}\]}
où $\underline{C}$ est la sous-variété linéaire canonique de $P_{\underline{T}}$ . Alors ci-dessus
$\overset{\hand{\sim}}{X}(C_{\hand{\unsure{\eta|y|\zeta}}}) \hand{\longrightarrow} \hand{Q}(C_{\hand{\unsure{\eta|y|\zeta}}})$ n'est autre que le morphisme des fibres génériques pour
le $\underline{T}$-morphisme $\overset{\hand{\sim}}{X}(\underline{C}) \hand{\longrightarrow} \hand{Q}(\underline{C})$ de ce dernier diagramme, et toute propri\'e-%
té constructible pour \var{\xout{\illegible{1}}}{}{typewriter strike-out}{ega25_25_091-v0005} le morphisme des fibres génériques implique
la propriété pour les fibres voisines. - Du point de vue notationel,
$\hand{Q}$ doit être envisagé (et même introduit) comme le nom du morphisme \var{\inserted[typed]{de foncteurs}}{de foncteurs}{insertion}{ega25_25_091-v0006}
naturel $\mathrm{Grass}^{m+1}(P) \hand{\longrightarrow}$ \var{\xout{Hilb}}{}{typewriter strike-out}{ega25_25_091-v0007} Sous-préschémas $(\mathrm{Grass}^{m}(P))$ .

15. \uline{Axiomatisation de certains des résultats précédents}.

Je pense surtout aux résultats des n°s 2 à 8, qui pour la plupart

\mspage{91}
